\subsection*{Note (5)}
The fifth note, on p.~55, adds an exception to the account in \S4,
p.~24, of preliminary work and geometric housekeeping. It identifies
the more delicate combinatorial interpretation of regular maps attached
to congruence subgroups and the effort to express their arithmetic
Galois action through the cyclotomic character. The ``cartes de
congruence'' thus return to both the regular maps of \S4 and the
cartographic arithmetic of \S3. The concrete change of perspective is
the one recorded in the final sentence: the combinatorial
``bi-icosaèdre'', first developed in a C4 course for geometric reasons,
subsequently provided a way to express the action on the category of
icosahedral maps, called here ``cartes de congruence d'indice 5''.
That wording is retained. The note also recalls, without developing it
further, the earlier work on cubic complexes; its emphasis is the
arithmetic use that emerged from a teaching construction.
(\href{https://grothendieck.umontpellier.fr/web/viewer.html?file=../119.pdf\#page=58}{Copy 1, p.~55, site 57};
\href{https://grothendieck.umontpellier.fr/web/viewer.html?file=../119.pdf\#page=27}{call in \S4, p.~24, site 26}.)
